\documentclass{tufte-handout}
\title{Notas de aula}
\author{Leonardo Suarez}
\date{\today}
\usepackage{analise_real}
\begin{document}
\maketitle
\section{Introdução}

\begin{defn}
    Seja \(A\) um conjunto. Dizemos que \(A\) é \emph{finito} se existe uma bijeção \(f : A \to \{1, 2, \ldots, n\}\) para algum \(n \in \N\); caso contrário, dizemos que \(A\) é \emph{infinito}.
\end{defn}

\bibliographystyle{plain}
\bibliography{referencias}
\end{document}
